\documentclass[aspectratio=169,11pt]{beamer}
\usepackage{amsmath,booktabs,array}
\definecolor{navy}{RGB}{30,65,98}
\definecolor{redproof}{RGB}{160,35,35}
\setbeamercolor{frametitle}{fg=navy}
\setbeamercolor{structure}{fg=navy}
\setbeamercolor{alerted text}{fg=redproof}
\setbeamertemplate{navigation symbols}{}
\setbeamertemplate{footline}{\hfill\usebeamercolor[fg]{structure}\insertframenumber\,/\,\inserttotalframenumber\hspace{5mm}\vspace{3mm}}
\setbeamersize{text margin left=9mm,text margin right=9mm}
\newcommand{\code}[1]{\texttt{#1}}
\newcommand{\cv}[1]{\mathtt{#1}}
\title{HMS target-geometry inconsistency}
\author{}
\date{}
\begin{document}

\begin{frame}{Code variables: the geometry inputs}
\small
\renewcommand{\arraystretch}{1.35}
\begin{tabular}{@{}p{.23\textwidth}p{.72\textwidth}@{}}
\toprule
Code variable & Physical meaning \\
\midrule
\code{reacty} & Laboratory vertical reaction position, read from \code{H.react.y}. Positive upward. \\
\code{xmis} & HMS vertical mispointing translation. The code places the vertex at \( -\cv{reacty}-\cv{xmis}\) in transport x, positive downward. \\
\code{zfoil} & Assigned foil position along the laboratory beam direction. \\
\code{c} & \(\cos(\cv{angle})\), with \code{angle} the spectrometer angle in radians. The code uses \(\cv{zfoil}\,\cv{c}\) as the longitudinal vertex position. \\
\code{L} & Distance along the spectrometer axis from its target reference plane to the sieve plane. HMS: 168 cm. \\
\code{xsT} & Nominal transport-x coordinate of the assigned sieve-hole center. Fixed by the hole grid. \\
\bottomrule
\end{tabular}
\vspace{2mm}

All positions and distances here are in cm. The target reference plane is at spectrometer z = 0.
\end{frame}

\begin{frame}{Code variables: the constructed fit targets}
\renewcommand{\arraystretch}{1.45}
\begin{tabular}{@{}p{.23\textwidth}p{.72\textwidth}@{}}
\toprule
Code variable & Physical meaning \\
\midrule
\code{xptarT} & Desired transport-x slope for the optics fit, \(dx/dz\). Dimensionless, usually reported in mrad. \\
\code{xtarT} & Corresponding transport-x intercept at the target reference plane, in cm. It need not equal the position at the foil. \\
\bottomrule
\end{tabular}
\vspace{4mm}

These are geometry-derived target values supplied to the fit.

\vspace{3mm}
\small
\textbf{Exact code correspondence.} In \code{targetTruth}, the input \code{xs} is the nominal hole coordinate stored as \code{xsT}. The outputs \code{t.xptar} and \code{t.xtar} become \code{xptarT} and \code{xtarT} in the fit tree.

\vspace{3mm}
The equations below consistently use \code{xsT}, \code{xptarT}, and \code{xtarT}.
\end{frame}

\begin{frame}{The condition to be satisfied}
\textbf{Show that:} the constructed slope and intercept describe a straight ray through the assigned sieve-hole center.

\vspace{4mm}
Starting at the target reference plane, propagation to the sieve adds distance times slope:
\[
\underbrace{\cv{xtarT}}_{\text{starting position}}
+\underbrace{\cv{L}\,\cv{xptarT}}_{\text{change in position}}
=\underbrace{\cv{xsT}}_{\text{hole center}}.
\]
\vspace{4mm}
Thus the required identity is
\[
\boxed{\cv{xtarT}+\cv{L}\,\cv{xptarT}-\cv{xsT}=0.}
\]
\small
All three positions must use the same transport coordinate origin. This is the reference ray to the nominal center, before considering finite hole size or scattering.
\end{frame}

\begin{frame}{The HMS equations implemented in the code}
The HMS branch of \code{targetTruth} constructs
\[
\cv{xptarT}=\frac{\cv{xsT}}{\cv{L}-\cv{zfoil}\,\cv{c}},
\]
\[
\cv{xtarT}=-\cv{reacty}-\cv{xmis}
-\cv{xptarT}\,\cv{zfoil}\,\cv{c}.
\]
\vspace{3mm}
The slope uses the hole position divided by the foil-to-sieve distance.

\vspace{3mm}
The intercept also includes the displaced vertex, \(-\cv{reacty}-\cv{xmis}\).

\vfill
\scriptsize Source: \code{spectrometer\_config.h}, HMS branch of \code{targetTruth}.\\
Storage mapping: \code{make\_fit\_ntuple\_from\_gmm.C}.
\end{frame}

\begin{frame}{Direct substitution exposes the contradiction}
\small
Substitute the implemented intercept, then the implemented slope:
\begin{align*}
&\cv{xtarT}+\cv{L}\,\cv{xptarT}-\cv{xsT}\\[2mm]
&=-\cv{reacty}-\cv{xmis}
-\cv{xptarT}\,\cv{zfoil}\,\cv{c}
+\cv{L}\,\cv{xptarT}-\cv{xsT}\\[2mm]
&=-\cv{reacty}-\cv{xmis}
+\cv{xptarT}(\cv{L}-\cv{zfoil}\,\cv{c})-\cv{xsT}\\[2mm]
&=-\cv{reacty}-\cv{xmis}+\cv{xsT}-\cv{xsT}\\[2mm]
&=\boxed{-\cv{reacty}-\cv{xmis}.}
\end{align*}
\normalsize
\alert{The required result is zero. The implemented result is generally nonzero.}

\vspace{2mm}
The pair reaches \(\cv{xsT}-\cv{reacty}-\cv{xmis}\), rather than \code{xsT}. Equality holds only when \(\cv{reacty}+\cv{xmis}=0\).
\end{frame}

\begin{frame}{What the contradiction means}
For the same longitudinal approximation, satisfying the identity requires
\[
\boxed{\cv{xptarT}=\frac{\cv{xsT}+\cv{reacty}+\cv{xmis}}
{\cv{L}-\cv{zfoil}\,\cv{c}}.}
\]
Use that same \code{xptarT} in the intercept equation.

\vspace{3mm}
\textbf{Conceptually:} the numerator must be the displacement from the vertex to the hole. The present numerator includes only the hole position, although the intercept retains the vertex displacement.

\vspace{3mm}
\small
\textbf{Scope.} This establishes the inconsistency of the constructed target pair. Its effect on fitted or replayed resolution requires a separate data comparison.

\vspace{3mm}
\footnotesize
\textbf{Replay context confirmed by Christine.} The 2024 replay used the same matrix as NPS replay, checked with the afterburner closure test. All zero-order offsets were removed for that replay. The HCANA version remains unchanged. These facts do not change the substitution above.
\end{frame}
\end{document}
